\documentclass[a4paper,11pt,twoside,openany,fleqn]{memoir} % gebruik A4 papier, 11pt lettergrootte, en een apparte titelpagina. Een nieuw hoofdstuk begint op eender welke kant (links of rechts). Vergelijkingen komen op een kleine identatie van de linker kant, en de nummering staat rechts
% Template voor een LaTeX thesis op ESAT. 
% Auteur: Dirk Van Hertem: dirk.vanhertem@ieee.org
% Verander eraan wat je wil, en stuur opmerkingen door, dan houd ik de
% versie up to date
% lees README voor meer uitleg, of kijk naar http://homes.esat.kuleuven.be/~dvherten/esatthesis.html

%%%%%%%%%%%%%%%%%%%%%%%%%%%%%%%%%
% Maak de titelpagina
%%%%%%%%%%%%%%%%%%%%%%%%%%%%%%%%%

\newcommand{\NAAM}{Uw Naam}
\newcommand{\NAAMb}{}% ook titelpagina.tex aanpassen als je dit invult (% teken weghalen)
\newcommand{\TITEL}{Beste thesis ooit al geschreven\ldots}
\newcommand{\SUBTITEL}{\ldots en der zal er geen betere meer komen}
\newcommand{\PROMOTOR}{Prof.\,dr.\,Weet~Alles}
%\newcommand{\PROMOTORb}{Prof.\,dr.\,Weet~Ook~Alles}% ook titelpagina.tex aanpassen (% teken weghalen)
\newcommand{\RICHTING}{Burgerlijk werktuigkundig-elektrotechnisch ingenieur, richting
elektrotechniek, optie energie}
\newcommand{\BEGELEIDER}{ir.\,Naam van de assistent}
\newcommand{\JAAR}{2008---2009}
%\newcommand{\BEGELEIDERb}{ir.\,Naam van assistent nr2}% ook titelpagina.tex aanpassen (% teken weghalen)


%%%%%%%%%%%%%%%%%%%%%%%%%%%%%%%%%
% Header file
%%%%%%%%%%%%%%%%%%%%%%%%%%%%%%%%%
% filename: header_esat.tex
% auteur: Dirk Van Hertem (dirk.vanhertem@ieee.org)
% deze  file zorgt dat de juiste  opties gezet zijn. Een  aantal opties staan niet  default, maar ik
% raad je aan ze eens te bekijken, en indien gewenst te gebruiken.


\tightlists % weinig extra witruimte bij opsommingen (memoir optie)

\usepackage[english,dutch]{babel} % nederlandse tekstopmaak % english
                                % toegevoegd in versie 0.4

\usepackage{graphicx} % Gebruik figuren
\graphicspath{{figuren/}}
\usepackage{tabularx} % maak een tabel met een bepaalde breedte (gebruikt in symbolenlijst.tex)
\usepackage{microtype} % toegevoegd in versie 0.4
\usepackage{ifthen}% toegevoegd in versie 0.5
\usepackage{calc}% nieuw in 0.6b om te kunnen rekenen in deze header
% \newboolean{fourierfont} % nieuw in versie 0.6b, op aanraden van Benoit Willems, naar 2009 toe zou ik
%                          % fourier  standaard willen zetten, maar  voorlopig nog niet  om geen grote
%                          %   veranderingen   te   krijgen    vlak   voor   de   deadline.   Vervangt
% % \usepackage{fourier} in commentaar ==> werkt nog niet, voor de volgende keer...
% \setboolean{fourierfont}{true}

% %\ifthenelse{\boolean{fourierfont}}{%

%%%%%%%%%%%%%%%%%%%%%
%%% Als je fourier wilt, uncomment de volgende regio:
%%%%%%%%%%%%%%%%%%%%%
%\makeatletter
%\let\my@@font@warning\@font@warning
%\let\@font@warning\@font@info
%\makeatother%
%\usepackage{fourier}%
%\makeatletter
%\let\@font@warning\my@@font@warning
%\makeatother

%%%%%%%%%%%%%%%%%%%%%
% inputenc
%%%%%%%%%%%%%%%%%%%%
% inputenc is afhankelijk van uw systeem (en os). Waarschijnlijk maakt
% een van de volgende \usepackage[]{inputenc} dat alles goed
% opgeslagen kan worden. Waarschijnlijk heb je latin1 nodig. \'e werkt altijd...

%\usepackage[latin1]{inputenc} % zorgt dat je ook � � en dergelijke kan ingeven
%\usepackage[utf8]{inputenc} % zorgt dat je ook � � en dergelijke kan ingeven

%%%%%%%%%%%%%%%%%%%%
% andere mogelijk nuttige pakketten
%%%%%%%%%%%%%%%%%%%%

% \usepackage{subfig}   % Gebruik subfiguren
% \usepackage{color}    % Gebruik kleuren
% \usepackage{icomma}   % zet de spati�ring goed bij kommagetallen (Europeese notatie)
% \usepackage{amssymb,amsmath,amsfonts} % extra wiskundige symbolen en
% constructies
% \usepackage{eurosym}  % \euro geeft het euroteken
% \usepackage{cite}     % verbetert de weergave van citaties 
% \usepackage{listings} % om code toe te voegen in uw document
% \usepackage{booktabs} % toegevoegd in version 0.4


% De volgende opties bepalen hoe figuren op een blad geplaatst moeten worden. Lees het volgende document voor meer informatie: http://www.ctan.org/tex-archive/info/epslatex.pdf

\setcounter{topnumber}{4}
\setcounter{bottomnumber}{4}
\setcounter{totalnumber}{10}
\renewcommand{\textfraction}{0.15}
\renewcommand{\topfraction}{0.85}
\renewcommand{\bottomfraction}{0.70}
\renewcommand{\floatpagefraction}{0.66}

% Opmerking: het volgende  stuk definieert de pagina opmaak. Indien de  keuze energie genomen wordt,
% dan wordt een opmaak met anderhalve regelafstand  genomen. Ofdat dit zinvol is, dat laat ik in het
% midden, de mensen die twijfelen lezen best de volgende url:
% http://www.tex.ac.uk/cgi-bin/texfaq2html?label=linespace. 
% De opmaak van energie heeft ook iets minder witruimte aan de randen. 
% Veranderen van opmaak kan simpelweg door de  waarde van \energie (zie regel na deze commentaar) de
% waarde "true" of "false" te geven.
\newboolean{energie}
\setboolean{energie}{false}

\ifthenelse{\boolean{energie}}% versie 0.6
{% deze optie is ook versie 0.6
% zet de pagina op (zie memman.pdf p 66)
\settrimmedsize{297mm}{210mm}{*}%
\setlength{\trimtop}{0pt}%
\setlength{\trimedge}{\stockwidth}%
\addtolength{\trimedge}{-\paperwidth}%
\settypeblocksize{654pt}{455pt}{*}% Van 634 naar 694 pt in versie 0.7
\setulmarginsandblock{2cm+\onelineskip+1.5\onelineskip}{2.65cm}{*}% Aanpassing 0.7
\setlrmargins{*}{*}{0.9}% verhouding buitenkantmarge/rugmarge. Verzet
% van 1.5 naar 0.9 vanaf versie 0.4
\setmarginnotes{17pt}{51pt}{\onelineskip}%
\setheadfoot{\onelineskip}{2\onelineskip}%
\setheaderspaces{*}{1.5\onelineskip}{*}% Aangepast versie 0.7
%
%zorg ervoor dat TOC niet te breed wordt:
%\setpnumwidth{2.5em}
%\setrmarg{3em} 
%
\checkandfixthelayout%
\OnehalfSpacing%
}%
{%
% zet de pagina op (zie memman.pdf p 66)
\settrimmedsize{297mm}{210mm}{*}%
\setlength{\trimtop}{0pt}%
\setlength{\trimedge}{\stockwidth}%
\addtolength{\trimedge}{-\paperwidth}%
\settypeblocksize{634pt}{448.13pt}{*}%
\setulmargins{4cm}{*}{*}%
\setlrmargins{*}{*}{0.9}% verhouding buitenkantmarge/rugmarge. Verzet
% van 1.5 naar 0.9 vanaf versie 0.4
\setmarginnotes{17pt}{51pt}{\onelineskip}%
\setheadfoot{\onelineskip}{2\onelineskip}%
\setheaderspaces{*}{2\onelineskip}{*}%
%
%zorg ervoor dat TOC niet te breed wordt:
%\setpnumwidth{2.5em}
%\setrmarg{3em} 
%
\checkandfixthelayout%
}
\setlength{\parindent}{0pt}
\setlength{\parskip}{0.5\baselineskip}
\setlength{\cftbeforechapterskip}{0pt}

% Voor de nummeringsdiepte te veranderen: met level 4 nummer je ook subsecties
\setcounter{secnumdepth}{4} % default gemaakt in versie 0.4
%\setcounter{tocdepth}{4}
\usepackage[hypertexnames=false]{hyperref} % geeft links zodat je een
                                % interactief document krijgt (zeer
                                % interessant als je het ooit online
                                % zet) 

% % caption maar 0.8  \textwidth breed maken, en normaal lettertype (dit was reeds  zo, maar nu kun je
% % zien wat je kan veranderen indien nodig) (versie 0.7)
% \changecaptionwidth
% \captionwidth{0.8\textwidth}
% %\normalcaptionwidth
% \captiontitlefont{\normalsize} ==> werkt nog niet, later nog eens zien

\usepackage{memhfixc}  % anders krijg je een error (memoir samen met hyperref...)

%%% Local Variables: 
%%% mode: latex
%%% TeX-master: "eindwerk_template"
%%% End: 


\begin{document}
\selectlanguage{english}%indien de thesis in het engels wordt gemaakt
%\selectlanguage{dutch}% voor normale thesissen

\vspace*{-1cm}\hspace*{-1cm}
\begin{minipage}{0.5\textwidth}
\textsc{Faculteit}\\ \textsc{Ingenieurswetenschappen}\\
\\
\textsc{Departement}\\
\textsc{Elektrotechniek -- ESAT}
\end{minipage}\hfill%
\ifthenelse{\boolean{energie}}{%extra  stuk  uit  versie  0.5a,  om te  compenseren  voor  de  gekke
                               %plaatsing van de sedes (met dank aan Joost Busschaert). Tijdelijke fix
\begin{minipage}{0.2\textwidth}
 \begin{flushright}
 \includegraphics[width=2cm]{sedes}
 \setlength{\unitlength}{1mm}
 \begin{picture}(28.5,2)
  \put(0,0){\line(1,0){28.5}}
 \end{picture}
 KATHOLIEKE\\UNIVERSITEIT\\LEUVEN
 \end{flushright}
 \end{minipage}
}{\begin{minipage}{0.3\textwidth}
 \begin{flushright}
 \includegraphics[width=2cm]{sedes}
 \setlength{\unitlength}{1mm}
 \begin{picture}(28.5,2)
  \put(0,0){\line(1,0){28.5}}
 \end{picture}
 KATHOLIEKE\\UNIVERSITEIT\\LEUVEN
 \end{flushright}
 \end{minipage}
}

\vfill
\pagestyle{empty}

\begin{center}
\begin{Huge}
\textsf{\TITEL}\\[4mm]
\end{Huge}
\begin{LARGE}
  \textsf{\SUBTITEL}
\end{LARGE}
\end{center}
\vfill
\begin{tabular}{p{0.49\textwidth}p{0.5\textwidth}}
   & Eindwerk voorgedragen tot het behalen van het diploma van
   \RICHTING\\
   & \\
   & \begin{Large}\textbf{\NAAM}\end{Large}\\
%  & \begin{Large}\textbf{\NAAMb}\end{Large}\\
   & \\
   & Promotor:\\
   & \hspace{4mm}\begin{large}\PROMOTOR\end{large}\\
%  & \hspace{4mm}\begin{large}\PROMOTORb\end{large}\\
   & Dagelijkse begeleiding:\\
   & \hspace{4mm}\begin{large}\BEGELEIDER\end{large}\\
%  & \hspace{4mm}\begin{large}\BEGELEIDERb\end{large}\\
\end{tabular}
\\
\\
\\
\begin{center}
\JAAR
\end{center}

\cleardoublepage

%%% Local Variables: 
%%% mode: latex
%%% TeX-master: "eindwerk_template"
%%% End: 


\pagenumbering{roman}
\pagestyle{plain} %%%% hier toevoegen 
\begin{otherlanguage}{dutch}
\begin{small}\vspace*{\fill} \copyright\ Copyright K.U.Leuven\\
\\
Zonder voorafgaande schriftelijke toestemming van zowel de
promotor(en) als de auteur(s) is overnemen, kopi\"eren, gebruiken of
realiseren van deze uitgave of gedeelten ervan verboden. Voor
aanvragen tot of informatie i.v.m. het overnemen en/of gebruik
en/of realisatie van gedeelten uit deze publicatie, wend U tot de
K.U.Leuven, Departement Elektrotechniek -- ESAT, Kasteelpark
Arenberg 10, B-3001 Heverlee (Belgi\"e). Telefoon~\mbox{+32-16-32 11
30} \& Fax.~\mbox{+32-16-32 19 86} of via email: \url{info@esat.kuleuven.be}.

Voorafgaande schriftelijke toestemming van de promotor(en) is
eveneens vereist voor het aanwenden van de in dit afstudeerwerk
beschreven (originele) methoden, producten, schakelingen en
programma's voor industrieel of commercieel nut en voor de
inzending van deze publicatie ter deelname aan wetenschappelijke
prijzen of wedstrijden.
\end{small}
\end{otherlanguage}
\vspace*{3cm}

\begin{otherlanguage}{english}
\begin{small}
\copyright\ Copyright by K.U.Leuven\\
\\
Without written permission of the promotors and the authors it is
forbidden to reproduce or adapt in any form or by any means any
part of this publication. Requests for obtaining the right to
reproduce or utilize parts of this publication should be addressed
to K.U.Leuven, Departement Elektrotechniek -- ESAT, Kasteelpark
Arenberg 10,  B-3001 Heverlee  (Belgium). Tel.~\mbox{+32-16-32 11
30}   \& Fax.~\mbox{+32-16-32 19 86} or by email: \url{info@esat.kuleuven.be}.

A written permission of the promotor is also required to use the
methods, products, schematics and programs described in this work
for industrial or commercial use, and for submitting this
publication in scientific contests.
\end{small}
\end{otherlanguage}

% \selectlanguage{dutch}

%%% Local Variables: 
%%% mode: latex
%%% TeX-master: "eindwerk_template"
%%% End: 
 % deze file zit al in de zip-file zitten

% include de volgende gedeelten per hoofdstuk. Werk best niet vanuit 1 file


\chapter*{Voorwoord}
\addcontentsline{toc}{chapter}{Voorwoord}
\label{cha:voorwoord}
Dank u iedereen. blablabla

%%% Local Variables: 
%%% mode: latex
%%% TeX-master: "eindwerk_template"
%%% End: 

\chapter*{Abstract}
\addcontentsline{toc}{chapter}{Abstract}
1 tot twee bladzijden over uw thesis

%%% Local Variables: 
%%% mode: latex
%%% TeX-master: "eindwerk_template"
%%% End: 

\tableofcontents\clearpage

\chapter*{Lijst van symbolen}
\addcontentsline{toc}{chapter}{Lijst van symbolen}

\begin{center}
  \begin{tabularx}{0.8\textwidth}{p{1.5cm}X}
    $\pi$ & het getal pi\\
    $42$  & The Answer to the Ultimate Question of Life, the Universe, and Everything\cite{h2g2}
  \end{tabularx}
\end{center}
%%% Local Variables: 
%%% mode: latex
%%% TeX-master: "eindwerk_template"
%%% End: 
\clearpage
\def\listfigurename{Lijst van figuren}
\listoffigures\clearpage
\def\listtablename{Lijst van tabellen}
\listoftables
%\def\lstlistlistingname{Lijst van codeblokken}

\cleardoublepage
\pagenumbering{arabic}
\pagestyle{ruled}


\chapter{Inleiding}
\label{cha:inleiding}

Inleidende tekst
\begin{figure}
  \centering
  \includegraphics[width=0.4\textwidth]{sedes}
  \caption{Dit is de sedes van  de kul. Der staat trouwens bij dat de KUL  in 1425 ontstaan is. Deze
    caption is  uitzonderlijk lang om te  testen ofdat het werkt  om de caption over  twee lijnen te
    zetten op 0.8 tekstbreedte.}
  \label{fig:sedes}
\end{figure}

\begin{itemize}
\item in deze lijst
\item staat een verwijzing
\item namelijk naar de figuur~\ref{fig:sedes}
\end{itemize}

%%% Local Variables: 
%%% mode: latex
%%% TeX-master: "eindwerk_template"
%%% End: 

\chapter{Het echte werk}
\label{cha:het-echte-werk}

\section{Titel om hier iets te hebben}
\label{sec:titel-om-hier}


Een proefschrift~\cite{wiki}, ook wel dissertatie of thesis genoemd, is een boek
geschreven door een promovendus met daarin een wetenschappelijke
verhandeling over een bepaald onderwerp. Het kan bestaan uit een
samenvoeging van al dan niet reeds eerder verschenen wetenschappelijke
publicaties. De promovendus doet hiermee verslag van een door hem of
haar zelfstandig beoefend wetenschappelijk onderzoek. Bij goedkeuring
van het proefschrift door zijn promotor krijgt de promovendus na een
succesvolle verdediging tijdens de promotieplechtigheid de titel van
doctor.

Het verrichten van dit onderzoek en het schrijven van een proefschrift
is in Nederland geen geringe opgave, waar men meestal een aantal jaren
voor uittrekt, niet zelden 4 of meer jaar. Om een proefschrift te
mogen schrijven moet men doorgaans eerst het doctoraalexamen of een
daarmee overeenkomstige titel hebben behaald, dus doctorandus (drs.),
meester in de rechten (mr.), of ingenieur (ir.) zijn. Sinds een aantal
jaren is het ook mogelijk om, bij grote uitzondering, zonder titel uit
het wetenschappelijk onderwijs te promoveren. Een persoon dient
hiervoor aannemelijk te maken dat hij in staat is zelfstandig
wetenschappelijk onderzoek te verrichten en met gerede kans op succes
een proefschrift kan voltooien. Hij dient hiervoor
wetenschapsbeoefenaren en een promotor te vinden die (schriftelijk)
garant voor hem staan. Echter, het is daarbij niet altijd even
gemakkelijk om in dat geval een promotor te vinden.

Na een promotie van een drs. komt de doctorstitel dr. voor de eerdere
titels te staan. De titel drs. verdwijnt dan. Bij mr. en ir. blijven
de titels staan. Indien het een promotie in de rechten betreft wordt
de titel mr. soms voor die van dr. geplaatst (mr.dr.). Veel juristen
blijven ook na hun promotie alleen de titel mr. voeren. Bij overige
promoties wordt dr. als eerste titel gehanteerd.

\section{Geschiedenis van het proefschrift}
\label{sec:geschiedenis-van-het}


Tot in de 20e eeuw was het niet ongebruikelijk om op stellingen te
promoveren. Er was in die gevallen geen promotieonderzoek gedaan. Vaak
promoveerde men op de dag waarop men was afgestudeerd.

\section{Verwante onderwerpen}
\label{sec:verwante-onderwerpen}

\begin{itemize}
\item Promotie~\cite{h2g2}
\item Wetenschappelijke promotie
\item Iuspromovendi (Promotierecht)
\item Thesis
\item Hoogleraar
\item Universiteit
\item scriptie\cite{pratchett06:_good_omens}
\item Titulatuur
\end{itemize}


Nu  gaan we  alles nog  eens herhalen  om beter  te zien  hoe het  eruit ziet  als  we verschillende
bladzijden hebben.


\section{Titel om hier iets te hebben}
\label{sec:titel-om-hier}


Een proefschrift~\cite{wiki}, ook wel dissertatie of thesis genoemd, is een boek
geschreven door een promovendus met daarin een wetenschappelijke
verhandeling over een bepaald onderwerp. Het kan bestaan uit een
samenvoeging van al dan niet reeds eerder verschenen wetenschappelijke
publicaties. De promovendus doet hiermee verslag van een door hem of
haar zelfstandig beoefend wetenschappelijk onderzoek. Bij goedkeuring
van het proefschrift door zijn promotor krijgt de promovendus na een
succesvolle verdediging tijdens de promotieplechtigheid de titel van
doctor.

Het verrichten van dit onderzoek en het schrijven van een proefschrift
is in Nederland geen geringe opgave, waar men meestal een aantal jaren
voor uittrekt, niet zelden 4 of meer jaar. Om een proefschrift te
mogen schrijven moet men doorgaans eerst het doctoraalexamen of een
daarmee overeenkomstige titel hebben behaald, dus doctorandus (drs.),
meester in de rechten (mr.), of ingenieur (ir.) zijn. Sinds een aantal
jaren is het ook mogelijk om, bij grote uitzondering, zonder titel uit
het wetenschappelijk onderwijs te promoveren. Een persoon dient
hiervoor aannemelijk te maken dat hij in staat is zelfstandig
wetenschappelijk onderzoek te verrichten en met gerede kans op succes
een proefschrift kan voltooien. Hij dient hiervoor
wetenschapsbeoefenaren en een promotor te vinden die (schriftelijk)
garant voor hem staan. Echter, het is daarbij niet altijd even
gemakkelijk om in dat geval een promotor te vinden.

Na een promotie van een drs. komt de doctorstitel dr. voor de eerdere
titels te staan. De titel drs. verdwijnt dan. Bij mr. en ir. blijven
de titels staan. Indien het een promotie in de rechten betreft wordt
de titel mr. soms voor die van dr. geplaatst (mr.dr.). Veel juristen
blijven ook na hun promotie alleen de titel mr. voeren. Bij overige
promoties wordt dr. als eerste titel gehanteerd.

\section{Geschiedenis van het proefschrift}
\label{sec:geschiedenis-van-het}


Tot in de 20e eeuw was het niet ongebruikelijk om op stellingen te
promoveren. Er was in die gevallen geen promotieonderzoek gedaan. Vaak
promoveerde men op de dag waarop men was afgestudeerd.

\section{Verwante onderwerpen}
\label{sec:verwante-onderwerpen}

\begin{itemize}
\item Promotie~\cite{h2g2}
\item Wetenschappelijke promotie
\item Iuspromovendi (Promotierecht)
\item Thesis
\item Hoogleraar
\item Universiteit
\item scriptie\cite{pratchett06:_good_omens}
\item Titulatuur
\end{itemize}

\cite{palaz2015analysis}
\cite{BiLSTM}

%%% Local Variables: 
%%% mode: latex
%%% TeX-master: "eindwerk_template"
%%% End: 

%\chapter{Besluit}
\label{besluit}
De masterproeftekst wordt afgesloten met een hoofdstuk waarin alle
besluiten nog eens samengevat worden. Dit is ook de plaats voor suggesties
naar het verder gebruik van de resultaten, zowel industri"ele toepassingen
als verder onderzoek.

\lipsum[1-7]

%%% Local Variables: 
%%% mode: latex
%%% TeX-master: "masterproef"
%%% End: 



% \nocite{eenboek} % refereert naar eenboek, zonder dat er in de tekst naar verwezen wordt
% \nocite{*} % verwijst naar alles in uw bibfile zonder dat er naar verwezen moet worden in de tekst

\bibliographystyle{IEEEtran_nl}   %    citeren   zoals    door   ieee   aangeraden    wordt.   Leest
                                % IEEEtran_bst_HOWTO.pdf voor  meer inlichtingen. IEEEtran_nl is een
                                % aangepaste versie van IEEEtran, waarbij de bindwoorden, maanden en
                                % dergelijke in het nederlands gezet zijn.
%\bibliographystyle{IEEEtran} % originele versie, te gebruiken als je uw tekst in het engels doet 
\bibliography{referenties} %  maak een eindwerk.bib bestand aan met uw referenties (dit is iets waar je best direct mee begint en niet op het einde doet)


%\appendix % indien gewenst
%\include{appendix1}

\end{document}

%%% Local Variables: 
%%% mode: latex
%%% TeX-master: t
%%% End: 
